%% kapa: exam=1 round=30 year=2019 period=1 form=A subject=economics subject_ko=경제학원론 seq=8 number=48
%% source: Q-Net 감정평가사 기출문제 (cst00309) artlSeq=5205205 file=extracted/1008/5205205/2019년 제30회 감정평가사 1차시험 문제지/2019년 제30회 감정평가사 1차시험 문제지/2019년도 제30회 감정평가사 1차 1교시 A형.hwp
%% conversion: kordoc 4.24.51 -> tools/gen1.py (auto); review=proofread-sonnet; flags=none
\begin{question}{48}
단기 비용곡선에 관한 설명으로 옳은 것을 모두 고른 것은? (단, 양(+)의 고정비용과 가변비용이 소요된다.)
\begin{framedtext}
ㄱ. 평균비용은 총비용곡선 위의 각 점에서의 기울기다. \par
ㄴ. 한계비용곡선은 고정비용 수준에 영향을 받지 않는다. \par
ㄷ. 생산량이 증가함에 따라 평균비용과 평균가변비용 곡선간의 차이는 커진다. \par
ㄹ. 생산량이 증가함에 따라 평균비용이 증가할 때 평균가변비용도 증가한다.
\end{framedtext}
\begin{choices}
  \item ㄱ, ㄴ
  \item ㄱ, ㄹ
  \item ㄴ, ㄷ
  \item ㄴ, ㄹ
  \item ㄷ, ㄹ
\end{choices}
\end{question}
\answer{4}
